\documentclass[11pt]{article}
\usepackage[margin=1in]{geometry}
\usepackage{amsmath,amssymb,amsthm,bm}
\usepackage{booktabs}
\usepackage{hyperref}
\usepackage{xcolor}
\newtheorem{lemma}{Lemma}
\newtheorem{proposition}{Proposition}
\newtheorem{remark}{Remark}
\theoremstyle{definition}
\newtheorem{definition}{Definition}
\newcommand{\R}{\mathbb{R}}
\newcommand{\lv}{\lambda_v}
\newcommand{\lr}{\lambda_r}
\newcommand{\lm}{\lambda_m}
\newcommand{\lrv}{\lambda_{rv}}
\newcommand{\Phixl}{\Phi_{x\lambda}}
\newcommand{\tf}{t_f}
\newcommand{\gap}[1]{\textcolor{red!70!black}{\textbf{[#1]}}}

\title{Second-order audit of the min-time catalog test:\\
what the free-time quotiented Jacobi determinant actually tests,\\
and what remains between it and a sufficient certificate}
\author{M.~Casey (host), prompted by GPT-6 Astra review \#2}
\date{2026-09-06}

\begin{document}
\maketitle

\begin{abstract}
The costate catalogs (18,360 min-time entries) carry a conjugate-point verdict
from \texttt{ms\_conjugate\_test} in its free-final-time form: the sign of
$\det[\Phixl(t,0)P,\; f(x(t))]$ sampled at the multiple-shooting junctions,
where $P$ spans the orthogonal complement of the converged initial costate
within the $(\lr,\lv)$ block. This note (i) states the problem with its actual
endpoint conditions (mass free at $\tf$, $\tf$ free), (ii) proves that the
instrument's quotient is equivalent to the Bonnard--Caillau--Tr\'elat (BCT)
free-time exponential-map construction on the $H=0$ slice, (iii) shows that the
free terminal mass adds no competitors in the all-burn problem, so the
point-target theorem applies after a mass reduction, and (iv) lists, without
euphemism, the hypotheses a per-entry \emph{sufficient} certificate would still
need and which of them are gates we can add. The honest description of a
current PASS is given in \S\ref{sec:meaning}.
\end{abstract}

\section{The problem as solved}\label{sec:problem}

State $x=(r,v,m)\in\R^7$ in the Earth--Moon rotating frame (ND units), control
$u=(\alpha,s)$ with $\alpha\in S^2$ and $s\in[0,1]$:
\begin{equation}\label{eq:dyn}
\dot r=v,\qquad \dot v=g(r)+h(v)+\frac{sT}{m}\alpha,\qquad \dot m=-\frac{sT}{c},
\end{equation}
with $g$ the gravity-plus-centrifugal field and $h(v)=(2v_y,-2v_x,0)$ the
Coriolis term. Minimize $\tf$ subject to $r(0),v(0),m(0)=m_0$ fixed and
$r(\tf)=r_f$, $v(\tf)=v_f$; \emph{$m(\tf)$ free, $\tf$ free}.

\paragraph{Normalized PMP.} With $\lambda_0=1$ and $\lambda=(\lr,\lv,\lm)$,
\begin{equation}\label{eq:H}
H(x,\lambda,u)=1+\lr\!\cdot v+\lv\!\cdot\big(g+h\big)+sT\Big(\frac{\lv\!\cdot\alpha}{m}-\frac{\lm}{c}\Big).
\end{equation}
Minimization over $\alpha$ gives $\alpha^*=-\lv/|\lv|$ (requires $|\lv|\neq0$);
then the coefficient of $s$ is $-T\,Q_{mt}$ with
\begin{equation}\label{eq:Qmt}
Q_{mt}:=\frac{|\lv|}{m}+\frac{\lm}{c},
\end{equation}
so $s^*=1$ wherever $Q_{mt}>0$. The catalog pipeline (pumpkyn \texttt{tfMin},
\texttt{ms\_tfmin}) \emph{assumes} $s\equiv1$ (``continuous burn is a
theorem'', \texttt{costate\_library\_methodology}); the sign condition
\eqref{eq:Qmt} is what that theorem asserts along an extremal, and it is
checkable from the junction states. Terminal conditions:
$r(\tf)=r_f$, $v(\tf)=v_f$, transversality $\lm(\tf)=0$ (free mass), and
$H(\tf)=0$ (free time; $H\equiv0$ since the system is autonomous).
Costate equations, with $s=1$:
\begin{equation}\label{eq:costate}
\dot\lr=-\big(\partial_r g\big)^{\!\top}\lv,\qquad
\dot\lv=-\lr-\big(\partial_v h\big)^{\!\top}\lv,\qquad
\dot\lm=-\frac{T|\lv|}{m^2}.
% SIGN CORRECTED 2026-09-10 (Astra proof review). H carries -T|lv|/m, so
% dH/dm = +T|lv|/m^2 and lam_m-dot = -dH/dm is NEGATIVE. The printed positive
% sign would have broken H = const. The CODE was always right: checked
% against the flown trajectory to ten digits, and max|H| along the arc is
% 3.3e-8. Document-only error -- but it is the document I reason from.
\end{equation}

\section{Two exact reductions}\label{sec:reductions}

\subsection{Mass is not a free terminal coordinate in the all-burn problem}
For every admissible all-burn trajectory, $m(t)=m_0-Tt/c$ regardless of
$\alpha(\cdot)$. Hence the set of competitors ``reaching $(r_f,v_f)$ at some
$\tf$ with $m(\tf)$ free'' is \emph{identical} to the set of 6-state trajectories
of
\begin{equation}\label{eq:dyn6}
\dot r=v,\qquad \dot v=g(r)+h(v)+\frac{T}{m_0-Tt/c}\,\alpha,
\end{equation}
reaching the \emph{point} $(r_f,v_f)$. The target manifold
$M_f=\{(r_f,v_f)\}\times\R_{>0}$ of the 7-state problem is a line whose free
direction is never exercised by an admissible variation
($\delta m(\tf)=-(T/c)\,\delta\tf$ is determined by $\delta\tf$). Consequently
\emph{the point-target, free-time second-order theory applies to
\eqref{eq:dyn6}}; no focal-point (manifold) construction is needed for the
min-time catalogs. (This is the opposite of the fixed-$\tf$ min-fuel line, where
the throttle varies and $m(\tf)$ is genuinely free; see \S\ref{sec:fixedtf}.)

\paragraph{What actually justifies the reduction (corrected 2026-09-10).}
The paragraph above argues from equality of \emph{competitor sets}, and that
argument is wrong: the stated problem admits a throttle $s\in[0,1]$, so the
all-burn trajectories are a strict subset and a first-order variation really
does produce extra mass freedom,
$\delta m_f=-k\,\delta\tf-k\!\int_0^{\tf}\!\delta s$ (Astra review,
2026-09-10). What H3 gives is the reduction \emph{at the level of critical
directions}: under a strict bang throttle the first-order variation
$\int_0^{\tf}-TQ_{mt}\,\delta s\,dt$ is non-negative for every feasible
$\delta s\le0$ and vanishes only for $\delta s=0$, so critical directions carry
no throttle component, and with no switch there is no switching-time variable.
That is \emph{not} by itself a second-order sufficiency theorem for the
throttled problem (\S\ref{sec:gaps}(a); wording corrected 2026-09-19).
\textbf{H3 is what we gate on.}

With the throttle explicit the Hamiltonian is \emph{affine} in $s$:
\[
  H(s)\;=\;1+\lr\!\cdot\!v+\lv\!\cdot\!(g+h)\;-\;s\,T\Big(\frac{|\lv|}{m}+\frac{\lm}{c}\Big)
       \;=\;H(0)-s\,T\,Q_{mt},
\]
so $Q_{mt}>0$ \emph{strictly} makes $s^{*}=1$ the unique minimiser with a
strict margin at every time: the throttle is \textbf{strictly bang}, and a
strictly bang component can be eliminated, leaving a smooth problem in the
direction alone --- which is the problem BCT addresses. Measured on the anchor:
$Q_{mt}\in[3.82,\,16.75]$, $dH/ds\le-0.672$ everywhere, and $H$ is affine in
$s$ to six digits ($H(0)=0.8922$, $H(\tfrac12)=0.4461$, $H(1)=0$).

\paragraph{An exact identity, and a free gate.} Since
$\frac{d}{dt}(\lm m)=\dot\lm m+\lm\dot m=-\frac{T|\lv|}{m}-\frac{T\lm}{c}=-T\,Q_{mt}$,
integrating and using $\lm(\tf)=0$, $m(0)=1$ gives
\[
  \boxed{\;\int_0^{\tf} T\,Q_{mt}\,dt \;=\; \lm(0)\;}
\]
--- the total strict-bang margin equals the initial mass costate. Verified on
the anchor: $5.5004$ against $\lm(0)=5.500404$. This is a cheap, exact
consistency check linking H3, the transversality condition and the mass
costate, and it is worth adding to the gate set: it uses only quantities the
flight already produces and it fails if any of the three is wrong.

The reduced system \eqref{eq:dyn6} is non-autonomous through $m(t)$. The
7-state autonomous form is its standard autonomization, but the two
Hamiltonians are \emph{not} equal (corrected 2026-09-19, external review; the
sentence that stood here said they agree). With $k=T/c$ the mass is a clock,
$m=m_0-kt$; the clock's costate is $p_\tau=-k\lm$ (indeed
$\dot p_\tau=-k\dot\lambda_m=k\,\partial_mH_6=-\partial_tH_6$), and
\[
  H_7\;=\;H_6+p_\tau\;=\;H_6-k\lm ,\qquad H_6=1+\lambda_{rv}\!\cdot\!f_{rv}.
\]
So $H_7\equiv0$ is a first integral, while $H_6=k\lm(t)$ varies along the arc
and vanishes only at $\tf$, where $\lm(\tf)=0$. The two \emph{formulations} are
equivalent under this transformation; the two Hamiltonians are not identical.
Nothing below depends on which form is used, but the instrument's row and column
choices are the 6-state ones, and \S\ref{sec:equiv} shows why that is right.

\subsection{Two exact invariances of the trajectory map}
Let $\mathcal T:\lambda(0)\mapsto x(\cdot;\lambda(0))$ be the map from initial
costate to the state trajectory of the PMP flow with $s\equiv1$,
$\alpha^*=-\lv/|\lv|$.
\begin{enumerate}
\item \textbf{Scale.} $\alpha^*$ depends on $\lv$ only through its direction,
and \eqref{eq:costate} is linear in $\lambda$; so for $a>0$,
$\mathcal T(a\lambda(0))=\mathcal T(\lambda(0))$. Differentiating,
\begin{equation}\label{eq:scale}
\Phixl(t,0)\,\lambda(0)=0\qquad\text{for all }t,
\end{equation}
where $\Phixl=\partial x(t)/\partial\lambda(0)$ (rows $r,v,m$; columns
$\lr,\lv,\lm$).
\item \textbf{Mass costate.} $\lm$ enters neither $\alpha^*$ nor the $(\lr,\lv)$
equations; so $\partial x(t)/\partial\lm(0)=0$: the $\lm$ column of $\Phixl$ is
identically zero, and $\partial m(t)/\partial\lambda(0)=0$ for every column.
\end{enumerate}
Both are exact, not approximate; they are why the na\"ive $7\times7$ or $6\times6$
determinant is identically singular (measured as $10^{-37}$--$10^{-13}$ noise
with spurious sign flips before the quotient was introduced, 2026-08-08).

\section{The BCT construction and the instrument's determinant}\label{sec:equiv}

\subsection{BCT free-time conjugate times}
For a normal extremal of a smooth control-affine system with a point target and
free final time, BCT (2007, \S3; also Agrachev--Sachkov Ch.~20--21) define the
exponential map
\[
\exp_{x_0}(t,\lambda)=\pi\circ e^{t\vec H}(x_0,\lambda),\qquad
\lambda\in\mathcal L:=\{\lambda:\ H(x_0,\lambda,u^*(\lambda))=0\},
\]
and call $t_c>0$ a \emph{conjugate time} if $(t,\lambda)\mapsto\exp_{x_0}(t,\lambda)$,
restricted to $\R\times\mathcal L$, is not an immersion at $(t_c,\lambda(0))$.
Its differential at $(t,\lambda(0))$ applied to $(\delta t,\delta\lambda)$ with
$\delta\lambda\in T_{\lambda(0)}\mathcal L$ is
\begin{equation}\label{eq:dexp}
d\exp\,(\delta t,\delta\lambda)=f(x(t))\,\delta t+\Phixl(t,0)\,\delta\lambda .
\end{equation}
Non-immersion means $\operatorname{rank}[\,\Phixl(t,0)|_{T\mathcal L},\ f(x(t))\,]<n$
with $n$ the state dimension.

\subsection{The tangent space of the $H=0$ slice}
At the minimizing control, $\partial H/\partial u=0$ tangentially on $S^2$ and
$s$ sits at its bound with $Q_{mt}>0$, so by the envelope argument
$\partial_\lambda\big[H(x_0,\lambda,u^*(\lambda))\big]=f(x_0,u^*)=:f_0$. Hence
\begin{equation}\label{eq:TL}
T_{\lambda(0)}\mathcal L=\{\delta\lambda\in\R^7:\ \delta\lambda\cdot f_0=0\},
\end{equation}
a 6-dimensional subspace. Note that pure $\lm$ shifts are \emph{not} tangent
($\delta\lambda\cdot f_0=\delta\lm\,\dot m_0$, $\dot m_0=-T/c\neq0$), and the
scaling direction $\lambda(0)$ itself is not tangent either
($\lambda(0)\cdot f_0=-1$ from $H=0$). So the $H=0$ slice is \emph{not} the
quotient by scaling; it is a transversal to it, which is the point of the
following lemma.

\subsection{Equivalence lemma}
Let $\lrv(0)\in\R^6$ denote the $(\lr,\lv)$ part of $\lambda(0)$, let $P\in\R^{6\times5}$
have orthonormal columns spanning $\lrv(0)^\perp$ (the instrument's
\texttt{null(qDir')}), and write $\Phi_{rv}$ for the $6\times6$ block of $\Phixl$
with rows $(r,v)$ and columns $(\lr,\lv)$.

\begin{lemma}\label{lem:image}
$\operatorname{Im}\big(\Phixl(t,0)|_{T\mathcal L}\big)\big|_{rv}
=\operatorname{Im}\big(\Phi_{rv}(t,0)\,P\big)$ for every $t$.
\end{lemma}
\begin{proof}
($\subseteq$) Take $\delta\lambda=(\delta\lrv,\delta\lm)\in T\mathcal L$. By the
$\lm$ invariance, $\Phixl\delta\lambda|_{rv}=\Phi_{rv}\delta\lrv$. Decompose
$\delta\lrv=a\,\lrv(0)+Pw$; by \eqref{eq:scale} restricted to the $rv$ rows
(the $\lm$ column being zero, \eqref{eq:scale} reads $\Phi_{rv}\lrv(0)=0$),
$\Phi_{rv}\delta\lrv=\Phi_{rv}Pw$.
($\supseteq$) Given $w\in\R^5$, set $\delta\lrv=Pw$ and choose
$\delta\lm=-\big(Pw\cdot f_{0,rv}\big)/\dot m_0$, which is possible because
$\dot m_0=-T/c\neq0$. Then $\delta\lambda\cdot f_0=0$, so
$\delta\lambda\in T\mathcal L$, and $\Phixl\delta\lambda|_{rv}=\Phi_{rv}Pw$.
\end{proof}

\begin{proposition}\label{prop:equiv}
For the reduced 6-state problem \eqref{eq:dyn6},
\[
\operatorname{rank}\big[\Phixl(t,0)|_{T\mathcal L},\ f(x(t))\big]_{rv}
=\operatorname{rank}\big[\Phi_{rv}(t,0)P,\ f_{rv}(x(t))\big],
\]
so $t$ is a BCT conjugate time iff the instrument's $6\times6$ determinant
$\det[\Phi_{rv}P,\,f_{rv}]$ vanishes. In particular the instrument's
\emph{free-time spec} (\texttt{stateRows=1:6, costateCols=8:13,
quotientDir=$\lambda(0)_{rv}$, freeTime=true}) is exactly the BCT immersion
test, not an ad hoc surrogate.
\end{proposition}
\begin{proof}
The flow column is common to both matrices; Lemma~\ref{lem:image} identifies
the column spaces of the other blocks; ranks of $[A,f]$ and $[B,f]$ agree when
$\operatorname{Im}A=\operatorname{Im}B$.
\end{proof}

\subsection*{The derivation this rests on, rebuilt (2026-09-10)}
\emph{The rank identity above is correct, but on its own it does not give the
proposition's ``so''. The construction as originally stated has a
\textbf{permanent kernel}, so it is non-immersive at every time, not merely at
conjugate times, and a rank drop cannot be read off it. This was found by an
adversarial review of this document (Astra, 2026-09-10) and is proved by our
own two lemmas. The derivation below replaces the earlier one.}

\paragraph{The permanent kernel.} Write $b=\dot m=-T/c\neq0$ and let $e_m$ be
the $\lm$ coordinate direction. Put
\[
  w \;=\; \lambda(0)+\tfrac1b\,e_m .
\]
Then $w$ is tangent to $\{H=0\}$, since
$dH_0(w)=f\!\cdot\!\lambda(0)+\tfrac1b f_m=-1+1=0$ (using $H=0\Rightarrow
\lambda\!\cdot\!f=-1$), and $w$ lies in the kernel of $\Phixl(t,0)$ for every
$t$, by \eqref{eq:scale} together with the vanishing $\lm$ column. Hence the
map on $\{H=0\}\times\R_t$ (tangent dimension $6+1=7$) has a one-dimensional
kernel at \emph{all} times. \emph{Verified numerically:}
$\|\Phixl(t,0)w\|/\|\Phixl\|\le5.5\times10^{-15}$ along the anchor arc.

\paragraph{The effective domain is six-dimensional.} Two facts make the
quotient concrete.
\begin{enumerate}
\item[(K1)] $T\{H=0\}$ projects \emph{isomorphically} onto the
  $(\lr,\lv)$ coordinates: $\delta\lambda\in T\{H=0\}$ means
  $\delta\lambda_{rv}\!\cdot\!f_{rv}+b\,\delta\lm=0$, so $\delta\lm$ is
  \emph{determined} by $\delta\lambda_{rv}$, which ranges over all of $\R^6$.
  Since the $\lm$ column of $\Phixl$ vanishes, the map factors through
  $\delta\lambda_{rv}$ alone.
\item[(K2)] $\Phi_{rv}(t)\,\lambda(0)_{rv}=0$ for every $t$ --- \eqref{eq:scale}
  with the vanishing $\lm$ column. \emph{Verified:}
  $\|\Phi_{rv}\lambda(0)_{rv}\|/\|\Phi_{rv}\|\le5.5\times10^{-15}$.
\end{enumerate}
So in the $\delta\lambda_{rv}$ coordinates the kernel $w$ becomes
$\lambda(0)_{rv}$, the effective domain is
$\R^6/\langle\lambda(0)_{rv}\rangle\oplus\R\,\partial_t$ of dimension
$5+1=6$, and the target $(r,v)(t)$ is six-dimensional. A conjugate time is
exactly a rank drop of this square induced map, whose matrix is
\[
  \big[\,\Phi_{rv}(t)\,P,\ f_{rv}(t)\,\big],
\]
with $P$ \emph{any} $6\times5$ basis of a complement of
$\langle\lambda(0)_{rv}\rangle$. That is the instrument's matrix.

\paragraph{The choice of complement does not matter.} Two complements differ
by a change of basis on the domain, so the determinants differ by a nonzero
constant, the same at every $t$; the vanishing set is therefore identical.
\emph{Verified} against a genuinely different (non-orthogonal) complement: the
ratio is constant to $6.9\times10^{-11}$ over the arc at $0.393694$, which is
precisely the reciprocal of the basis-change determinant $2.540047$, and both
choices report the same sign changes.

\begin{remark}[What this repair does and does not fix]
It establishes the \emph{instrument}: the $6\times6$ determinant is the rank
test of the correct six-dimensional induced map, derived rather than asserted.
It does \emph{not} by itself transfer the BCT sufficiency theorem --- see the
scope caveat on all-burn competitors in \S\ref{sec:hyp}.
\end{remark}

\begin{remark}[Why rows $1{:}6$ and not $1{:}7$]
In the 7-state form the mass row of $\Phixl$ is identically zero
(\S\ref{sec:reductions}), so a $7\times7$ block would be structurally singular
for a reason that has nothing to do with conjugacy. The reduction of
\S2.1 makes the 6-state block the correct object. This is the min-time
counterpart of the 2026-09-05 fixed-$\tf$ correction (rows $[1{:}6\ 14]\to1{:}7$),
and it goes the \emph{other} way: there the mass row carries information
(throttle varies), here it carries none.
\end{remark}

\begin{remark}[The quotient is not a symmetry of the BVP]
Astra's objection stands and is answered by the lemma rather than dismissed:
with $\lambda_0=1$ and $H(0)=0$, scaling $\lambda\mapsto a\lambda$ is not a
symmetry of the boundary-value problem (it violates $H=0$), so ``quotient by
scaling'' is not literally BCT's construction. What the lemma shows is that the
\emph{image} of the tangent space of the correct construction ($H=0$ slice)
coincides with the image of the scaling quotient, because the two differ only
by directions ($\lambda(0)$ itself, and $\lm$) that the trajectory map kills
exactly. The determinant is therefore the same function of $t$ up to a nonzero
factor, and its sign changes are the same events.
\end{remark}

\section{What a sufficient certificate needs}\label{sec:hyp}

The theorem, as BCT state it (checked against the paper, 2026-09-18; the
earlier ``Thm.~3.x'' was not an identifiable statement). Their standing
assumption is that the set $U$ of control values is \emph{open} in a manifold
$N$ (``the bang-bang case is not treated here''), with $N=S^2$, $U=N$ for orbit
transfer (their footnote~1). Under
\textbf{(L)} the strong Legendre condition along the extremal and
\textbf{(S)} \emph{strong regularity}: the control is of corank one on
\emph{every subinterval} $[t_1,t_2]\subset[0,T]$,
\textbf{Theorem~2.12} reads: \emph{the first geometric conjugate time coincides
with the first conjugate time $t_c$; the trajectory is locally optimal on
$[0,t_c)$ in the $L^\infty$ topology, and in the $C^0$ topology whenever the
extremal is normal; for $t>t_c$ it is not locally optimal in $L^\infty$ topology
on $[0,t]$.} (Theorem~1.13 gives the same sufficiency for a corank-one extremal.)
For a normal extremal with free final time the geometric conjugate times are
computed by their \emph{Test~3} (\S2.3, item (1)(b)): a zero of
$\det(\delta x_1,\dots,\delta x_{n-1},\,f)$ over the Jacobi fields tangent to
$\{H_r=0\}$ --- which is the instrument's determinant, \S\ref{sec:equiv}.
Two words of the sentence that stood here were not the theorem's: it said
``\emph{strict} strong local minimizer''; BCT conclude \emph{local optimality in
the $C^0$ topology}, without ``strict''. Applied to \eqref{eq:dyn6} via
\S\ref{sec:reductions}, each hypothesis becomes a per-entry statement:

\begin{center}\small
\begin{tabular}{@{}p{0.30\textwidth}p{0.36\textwidth}p{0.28\textwidth}@{}}
\toprule
hypothesis & what it means here & status / gate \\
\midrule
H1 normal extremal & a lift with $\lambda_0=1$ exists & \textbf{have} (the BVP is solved with $\lambda_0=1$ to $10^{-10}$; pumpkyn agrees to $|\Delta z|\sim10^{-9}$). Abnormal lifts of the same trajectory: \textbf{excluded by the $\dim S=1$ probe} (\S\ref{sec:abnormal}), built and passing on the golden cells \\
H2 strong Legendre & $H_{\alpha\alpha}|_{TS^2}=\tfrac{T}{m}|\lv|\,I\succ0$, i.e.\ $|\lv(t)|>0$ on $[0,\tf]$ & \textbf{gated with a margin, over the whole arc} (2026-09-18, \S\ref{sec:gaps}): a lower bound on $\min_{[0,\tf]}|\lv|$ must exceed $10^{-5}$. Golden cells, sampled: $\min|\lv|=0.091$ (dro), $0.012$ (halo), $0.143$ (dpo) \\
H3 all-burn is the PMP control & $Q_{mt}(t)>0$ on $[0,\tf]$, eq.~\eqref{eq:Qmt} & \textbf{gated with a margin, over the whole arc} (same bound and floor): $\min Q_{mt}=0.099$, $0.045$, $0.178$; this is \texttt{tfMinEoM}'s own switching function, so it also certifies that the event-driven propagator flew no switch \\
H4 smoothness / no collision & $m>0$, $r$ away from the primaries & \textbf{have} (\texttt{true\_min\_altitude} screen) \\
H5 no conjugate time on $(0,\tf]$ & $\det[\Phi_{rv}P,f_{rv}]\neq0$ on $(0,\tf]$ & \textbf{assessed}, not proved: sign changes at $K$ junctions (now including $\tf$); even-multiplicity zeros and pairs within one segment are invisible \\
\bottomrule
\end{tabular}
\end{center}

\subsection{What the theorem does not cover, and what is now checked (2026-09-18)}\label{sec:gaps}
Reading the source, and the two reviews of 2026-09-17 (FINDINGS 78, 80), leave
five items. Each is stated with what was done about it.

\paragraph{(a) The throttle is outside BCT's setting.}
BCT assume an open control set; theirs is the problem with the control on
$S^2$, which is ours \emph{after} the throttle is fixed at $s=1$
(\S\ref{sec:reductions}). Against competitors that throttle, $s\in[0,1]$ is a
closed set with the extremal on its boundary, and Theorem~2.12 does not speak.
The theory that does is the second-order sufficiency theory for controls with a
\emph{continuous} component entering nonlinearly and a \emph{bang-bang}
component entering linearly (Osmolovskii--Maurer 2006, 2009; their 2012
monograph; Maurer--Osmolovskii 2004 for the time-optimal bang-bang case). Our
bang-bang component has \emph{no switch}: under the \emph{strict bang-bang
property}, here $Q_{mt}>0$ on all of $[0,\tf]$, the critical cone contains no
switching-time variation and the quadratic form reduces to the smooth one in
the direction. \textbf{That reduction is cited, not instantiated here}: their
continuous component is unconstrained in $\mathbb{R}^m$ and ours lives on
$S^2$, so a chart argument is owed. What is exact and checkable is the gap
\[
 H(s,\alpha)-H(1,\alpha^{*})\;=\;T\,Q_{mt}\,(1-s)\;+\;\frac{s\,T\,|\lv|}{2m}\,\|\alpha-\alpha^{*}\|^{2},
\]
whose first coefficient is H3 and whose second is H2 \emph{times} $s$: H2 is
the positive angular Hessian at the reference throttle $s=1$, not a uniform
coefficient over all throttles (the second term vanishes at $s=0$, where the
first carries the whole gap). What the Osmolovskii--Maurer framework would
still require, beyond a chart on $S^2$: the applicable theorem named, its
endpoint qualifications and regularity, the second-variation (coercivity)
condition, the free-time treatment, and the topology of the asserted minimum.

\paragraph{(b) H2 and H3 were sampled positivity with no margin.}
The gate was $\min_k|\lv(t_k)|>0$ and $\min_k Q_{mt}(t_k)>0$: it passed
$10^{-300}$, passed values below the costates' own numerical error
($\sim10^{-6}$ in $\lm$ on the flown arc), and said nothing between samples.
Both now use one slope bound. In the CR3BP
$\dot\lambda_v=-\lr-C^{\!\top}\lv$ with $C=\partial_v f_v$ the Coriolis matrix,
which is \emph{skew}; so it drops out of
$\tfrac{d}{dt}|\lv|=\hat\lambda_v\!\cdot\!\dot\lambda_v=-\hat\lambda_v\!\cdot\!\lr$
and
\[
  \Big|\tfrac{d}{dt}|\lv|\Big|\;\le\;|\lr|
\]
(external review, 2026-09-19; the code's first version used the looser
$|\lr|+2|\lv|$). On an all-burn arc, with
$\dot m=-T/c$ and $\dot\lambda_m=-T|\lv|/m^{2}$, the two mass terms of
$\dot Q_{mt}$ cancel:
\[
 \dot Q_{mt}\;=\;\frac{1}{m}\,\frac{d|\lv|}{dt}\qquad\text{exactly.}
\]
If $L_k$ bounds $|v'|$ \emph{on the whole interval}, then
$\min_{[t_k,t_{k+1}]}v\ \ge\ \tfrac12(v_k+v_{k+1})-\tfrac12L_k\,\Delta t_k$
(\texttt{between\_sample\_bound}). \textbf{The code takes $L_k$ as the larger of
the slope bound at the interval's two ENDS, and from that definition the
inequality does not follow}: end samples cannot see an interior excursion
($v=1-a\sin^2(\pi t/h)$ has zero slope at both ends and dips to $1-a$). What is
computed is therefore a \emph{lower-bound estimate}, not a bound. A rigorous
version is cheap in principle: with $\|\partial_rf_v\|\le\kappa$ on the
interval, $(|\lr|,|\lv|)$ obeys the comparison system
$\dot a\le\kappa b,\ \dot b\le a$, so
$|\lr(t)|\le a_0\cosh(\sqrt\kappa h)+\sqrt\kappa\,b_0\sinh(\sqrt\kappa h)$, which
is a valid $L_k$ --- given a lower bound on the distances to the primaries over
the interval, which end samples again do not supply. Not implemented.
Note also that H2 \emph{implies} H3 on an exact all-burn extremal:
$\lm(t)=\int_t^{\tf}T|\lv|/m^2\,d\tau\ge0$, so $Q_{mt}\ge|\lv|/m>0$; H3 is a
numerical consistency check, not an independent hypothesis. The certifier and the catalog audit gate the
\emph{estimate} above $10^{-5}$. (That floor was motivated as ten times the
measured $\lm$ uncertainty; that is not an error budget for either quantity ---
$\epsilon_Q\le\epsilon_{\lambda_v}/m+\epsilon_{\lm}/c$, and the error in
$|\lv|$ is governed by the error in $\lv$. The observed margins are five orders
larger than the floor, so nothing turns on it here.) On the
$70$\,mN library ($576$ entries, $\sim7000$ samples per flight, steps
$\le0.0035$) the bound sits within $1\%$ of the sampled minimum; the smallest
bounds are $0.2785$ ($|\lv|$) and $0.3126$ ($Q_{mt}$). This is an estimate on a
finely sampled flight, not a validated enclosure: $L_k$ is taken at the
interval's ends.

\paragraph{(c) Strong regularity (S) is checked on the whole arc only.}
$\dim S=1$ (\S\ref{sec:abnormal}) is corank one on $[0,\tf]$. BCT's (S) asks
it on every subinterval; Theorem~2.12, which identifies the geometric
conjugate time we compute with the first conjugate time, is stated under (S).
\textbf{Not checked.} The lift-space rank on prefixes $[0,t]$ and on a few
interior windows would be additional numerical probes; they would not check a
condition quantified over \emph{every} subinterval. A proof would have to
propagate the rank property, e.g.\ by analyticity away from collisions and
$\lv=0$; that argument is not supplied.

\paragraph{(d) The certificate is for FIXED phases.}
The endpoints are points, $x(0)=x_D(s_D)$ and $x_{rv}(\tf)=x_A(s_A)$, and the
variation space of the conjugate test is that of the point-to-point problem.
Nothing above addresses sliding an endpoint along its orbit, which is the
question a phase catalog is consulted for; a saddle of the free-phase problem
passes by construction. The first-order part is free. With
$H=1+\lambda\!\cdot\!f$ and $\lambda\!\cdot\!f=-1$, $\lambda$ is the gradient
of the cost-to-go, and
\[
 \frac{\partial\tf}{\partial s_D}=+\lambda_{rv}(0)\cdot x_D'(s_D),\qquad
 \frac{\partial\tf}{\partial s_A}=-\lambda_{rv}(\tf)\cdot x_A'(s_A)
\]
(\texttt{phase\_sensitivity}; the mass costate does not enter: $m(0)$ does not
depend on the phase and $m(\tf)$ is free). Two uses.
\emph{Cross-check X3} (\texttt{phase\_transversality\_check}): re-solve the
transfer at $s\pm2\times10^{-4}$ and difference the two flight times. It is a
\emph{finite-difference consistency test}: the observable does not use the
returned costates, but the re-solve is an indirect solve seeded from the stored
ones in the same implementation, and each derivative tests one scalar
projection of an endpoint costate, not the whole vector. (The independent
oracle is the analytic single-integrator test of the formula.) Results: on the
hand-built library, twelve of twelve sampled derivatives agree, $4\times10^{-8}$
to $9\times10^{-4}$ relative; on the rebuilt library now of record, eleven
agree ($6\times10^{-8}$ to $4\times10^{-4}$) and one re-solve did not
converge, so the verdict there is UNRESOLVED, not PASS. The formulas are
smooth-branch sensitivities: they equal the gradient of the minimum-time value
only where that value is differentiable and this branch realises it.
\emph{Stationarity:} an entry is a first-order candidate for the
orbit-to-orbit minimum where both sensitivities vanish. None of the $576$
\emph{stored roots} is: the fastest, $16.226$\,d, has
$(\partial_{s_D},\partial_{s_A})\tf=(-0.92,-2.05)$ d per unit phase. That says
where to look next on this branch; it does not exclude an unrecorded branch
with a stationary root at a grid phase, nor locate the global minimum. The
second-order free-phase condition (endpoint tangent and curvature terms in the
accessory problem) is a separate problem and is not attempted.

\paragraph{(e) H5 remains assessed.}
Unchanged: candidate-triggered sampling with a policy floor; the uncovered
initial interval is reported, not bounded, and the short-time sign of the
determinant has not been derived. Two consumer defects were closed: a scan
reporting \texttt{multiplicity}${}>0$ with its other counts zero certified,
and the certifier's test seam could return a shippable certificate.

\subsection{On normality}\label{sec:abnormal}
An abnormal extremal is a solution of the same BVP with $\lambda_0=0$, i.e.\ with
the ``$1$'' removed from \eqref{eq:H} and $H\equiv0$ still imposed. Our accepted
lift has $\lambda_0=1$ and finite $\lambda$, which proves the trajectory
\emph{admits} a normal lift. BCT's theorem is stated for normal extremals in
this sense (existence of a normal lift under the strong Legendre condition, the
minimizing control being uniquely determined). What it does not address is
whether the same trajectory also admits an abnormal lift.

\emph{Corrected 2026-09-10 (Astra proof review).} The sentence that stood here
claimed an abnormal lift ``would require $\lv\equiv0$''. That is false:
abnormality means the \emph{cost} multiplier $\lambda_0$ vanishes, and an
abnormal extremal may perfectly well have $|\lv|>0$ and obey the same unique
direction law. The informal argument built on it is withdrawn; the computation
below stands on its own and is what we actually use. The lifts of a \emph{fixed} trajectory form a
linear space
\[
S=\{\lambda(\cdot):\ \dot\lambda=-A(t)^{\!\top}\lambda,\ \ \lv(t)\parallel\alpha(t)\ \forall t,\ \ \lm(\tf)=0\},
\]
with $A(t)=\partial f/\partial x$ at the \emph{frozen} flown control. Along any
$\lambda\in S$ the quantity $\lambda\cdot f$ is conserved (it is $H-1$ for the
normal lift, and $d(\lambda\cdot f)/dt=\lambda\cdot\partial_u f\,\dot u=0$ at a
stationary control), so $\lambda\mapsto\lambda\cdot f$ is a linear functional on
$S$ whose kernel contains the abnormal lifts. Our normal lift has
$\lambda\cdot f=-1\neq0$. Therefore \emph{$\dim S=1$ EXCLUDES an abnormal lift}
(if $\dim S\ge2$ the kernel is nontrivial). $\dim S=7-\operatorname{rank}C$ with
$C=\big[\,[\alpha_k]_\times\Psi_v(t_k)\ ;\ e_m^{\!\top}\Psi(\tf)\,\big]$ stacked over
samples $t_k$, $\Psi$ the fundamental matrix of the fixed-control adjoint.
Implemented in \texttt{mintime\_hypothesis\_gates}: on the three golden cells
$\dim S=1$ with singular-value ratios $\sigma_7/\sigma_6=1.8\times10^{-7}$,
$4.3\times10^{-8}$, $5.4\times10^{-8}$; the accepted $\lambda(0)$ lies in $S$
to $10^{-7}$ (integration tolerance) and $\lambda\cdot f=-1$ to $10^{-10}$ along
each arc. (The positivity of the multiplier $\rho$ in $\lv=-\rho\alpha$ is
automatic for elements of a one-dimensional $S$ containing the normal lift.)

\begin{remark}[Which direction we may claim]
$S$ imposes the linear adjoint and stationarity conditions and
$\lv\parallel\alpha$, but \emph{not} the PMP inequalities ($\lv=-\rho\alpha$
with $\rho\ge0$, and $Q_{mt}\ge0$). A nonzero kernel element therefore need not
be a genuine minimising multiplier. Consequently $\dim S=1$ together with a
normal element of nonzero functional value is a \textbf{sufficient algebraic
exclusion} of abnormal lifts --- which is the direction the gate is used in ---
and the converse is \emph{not} established. The earlier ``iff'' is withdrawn.
Making the converse available would mean testing the kernel against the PMP
cone, not merely counting its dimension.
\end{remark}

\subsection{On the sampled determinant}\label{sec:sampling}
Even with H1--H4 established, H5 is checked by sign changes at $K$ points.
Three failure modes remain: (a) two conjugate times inside one segment;
(b) a double root (even multiplicity, no sign change) --- for two decoupled
copies of a problem with a simple conjugate time this is generic, so it is not
exotic; (c) the endpoint case $\det(\tf)=0$, which the instrument now reports as
\texttt{ENDPOINT} rather than a refutation. Mitigations, in order of cost:
sample the determinant along each segment's variational integration (the STM
is already integrated densely; only the determinant is discarded);
monitor $\sigma_{\min}$ of the equilibrated block as a continuous quantity, so
an even-multiplicity zero shows as a dip to the noise floor; and, if a
\emph{rigorous} statement is wanted, validated integration of the variational
equations with a sign-definite enclosure of the determinant --- out of scope
for this program.

\subsection*{H6: excluding the reduced problem's spurious zeros (new, 2026-09-10)}
The reduction to \eqref{eq:dyn6} carries a hazard the autonomous form does not.
Writing $h(t)=p(t)\!\cdot\!F(t,y,u^{*})$ for the \emph{reduced} Hamiltonian and
$J(t)=D_{p_0}y(t)$, differentiating along the variational equations gives
$p(t)^{\top}J(t)=0$, so whenever $\operatorname{rank}J=5$,
\[
  \det[\,J P,\ F\,]=0 \iff h(t)=0 .
\]
In the autonomous normal point-target problem $h\equiv-1$ kills this mechanism.
In the reduction it does not, and H1--H4 do not exclude it: the determinant
could vanish for a reason that is \emph{not} a loss of optimality (Astra proof
review, 2026-09-10). Left there, this would undercut the instrument.

It is, however, computable in closed form here. From $H=1+\lambda\!\cdot\!f=0$,
\[
  \lambda_{rv}\!\cdot\!f_{rv}
  \;=\;-1-\lm\Big(-\frac{T}{c}\Big)\;=\;-1+\frac{T}{c}\,\lm ,
\]
which vanishes exactly at $\lm=c/T$. And $\lm$ is \emph{monotonically
decreasing} ($\dot\lm=-T|\lv|/m^{2}<0$) from $\lm(0)$ to $\lm(\tf)=0$. Hence

\begin{center}\fbox{\textbf{H6:} the spurious mechanism cannot fire iff
$\lm(0)<c/T$.}\end{center}

\noindent This is exact, costs nothing, and by the identity above is equivalent
to $\int_0^{\tf}T\,Q_{mt}\,dt<c/T$ --- the total strict-bang margin must stay
below $c/T$. \emph{Measured over the shipped 70\,mN catalog}: $c/T=49.38$,
$\lm(0)\in[2.11,\,10.88]$ across all 53 entries, so the margin is $38.50$ and
the headroom a factor of $4.5$; $|\lambda_{rv}\!\cdot\!f_{rv}|\ge0.78$
everywhere, and along the anchor arc it stays in $[-1.0000,-0.8886]$. \textbf{No
catalog entry can fire it.} \emph{Wiring note:} H6 is not yet computed by
\texttt{mintime\_hypothesis\_gates}; it should be, since it is one subtraction.

\section{What a PASS means today}\label{sec:meaning}
For a catalog entry with \texttt{conj\_pass = 1} the defensible statement is:
\begin{quote}
\emph{The entry is a normal PMP extremal of the all-burn min-time problem
(residual $10^{-10}$, independently reproduced by pumpkyn \texttt{tfMin}
to $|\Delta z|\lesssim10^{-9}$), and the Bonnard--Caillau--Tr\'elat free-time
conjugate determinant, in the exactly equivalent quotiented form of
Proposition~\ref{prop:equiv}, shows no sign change at the $K$ junction samples
on $(0,\tf]$ (an initial arc where the state block is structurally rank-deficient
having been skipped).}
\end{quote}
With \texttt{mintime\_hypothesis\_gates} (H1$'$, H2, H3 on the dense flight,
passing on the golden cells) the statement becomes: \emph{the entry is a
normal extremal admitting no abnormal lift, with the strong Legendre condition
and the all-burn condition holding with a margin over the whole flown arc, and
no sign change of the BCT conjugate determinant at the sampled times on
$(0,\tf]$.}
\emph{Revised 2026-09-18.} That is what is measured. It is \emph{not} every
hypothesis of Theorem~2.12: H5 is sampled rather than proved, (S) is checked on
the whole arc only, and the throttle puts the original problem outside BCT's
open-control setting (\S\ref{sec:gaps}). The sentence a methods paper can carry
is: \emph{each entry is a normal extremal of the fixed-phase all-burn problem
that passes a numerical assessment of the hypotheses of BCT's Theorem~2.12; if
those hypotheses hold it is locally optimal in the $C^0$ topology among
all-burn trajectories with the same endpoints.} Running the gates over the catalogs is a
one-hour campaign (one flight per entry) and is the natural next step.

\section{The fixed-$\tf$ line, for contrast}\label{sec:fixedtf}
In the smoothed-fuel problems the throttle varies, so $m(\tf)$ is genuinely free
and $M_f=\{(r_f,v_f)\}\times\R_{>0}$ is a real target manifold. The interior test
with rows $1{:}7$ (2026-09-05) is the fixed-endpoint Jacobi \emph{necessary}
condition: a conjugate point of the point-target accessory problem yields an
admissible zero-second-variation direction for the manifold problem too (a
subset of its competitors). Sufficiency for the manifold problem needs the focal
construction --- the Lagrangian subspace $\{\delta r=\delta v=0,\ \delta\lm=0\}$
at $\tf$ propagated \emph{backward} and tested for intersection with the vertical
--- i.e.\ the backward Riccati sweep of the certification survey, item~3.2.
Astra review \#2 also established that for the continuous clipped laws
(\texttt{eps}, \texttt{huberc}) the branchwise-AD STM is the true derivative of
the flow at transversal band crossings ($F^+=F^-$ there), so the instrument's
input is right; what is missing is the terminal-manifold construction, not the
STM.

\section{Recommended additions to the instrument (min-time)}
\begin{enumerate}
\item \textbf{Done:} \texttt{mintime\_hypothesis\_gates} gates $\min_t|\lv|$,
$\min_t Q_{mt}$ and $\dim S$ on the dense flight (\texttt{tests/test\_mintime\_gates},
15/15 on the golden cells). Next: store all three per catalog entry
(\texttt{conj\_catalog\_pass} already flies each entry).
\item Dense determinant sampling: in \texttt{collectSTMs}, keep the variational
ODE's dense output for each segment and evaluate the quotiented determinant at
every integrator step; report the finest bracket of each sign change and the
minimum of $\sigma_{\min}$.
\item \textbf{Done} on the golden cells (\S\ref{sec:abnormal}); catalog scale
with the item above.
\item Rename the catalog field's documentation from ``conjugate-certified'' to
the statement of \S\ref{sec:meaning}.
\end{enumerate}

\section*{References}
\begin{itemize}
\item B.~Bonnard, J.-B.~Caillau, E.~Tr\'elat, ``Second order optimality
conditions in the smooth case and applications in optimal control,''
\emph{ESAIM: COCV} 13(2) (2007) 207--236, doi:10.1051/cocv:2007012. Theorems 1.13 and 2.12;
the free-final-time Test~3 of \S2.3.
\item N.~P.~Osmolovskii, H.~Maurer, ``Second order optimality conditions for
controls with continuous and bang-bang components,'' in \emph{Systems, Control,
Modeling and Optimization}, IFIP, Springer, 2006, pp.~297--307,
doi:10.1007/0-387-33882-9\_28.
\item N.~P.~Osmolovskii, H.~Maurer, ``Second order sufficient optimality
conditions for a control problem with continuous and bang-bang control
components: Riccati approach,'' in \emph{System Modeling and Optimization}, IFIP
AICT, Springer, 2009, pp.~411--429, doi:10.1007/978-3-642-04802-9\_24.
\item N.~P.~Osmolovskii, H.~Maurer, \emph{Applications to Regular and Bang-Bang
Control}, SIAM, 2012, doi:10.1137/1.9781611972368.
\item H.~Maurer, N.~P.~Osmolovskii, ``Second order sufficient conditions for
time-optimal bang-bang control,'' \emph{SIAM J.\ Control Optim.} 42(6) (2004)
2239--2263, doi:10.1137/S0363012902402578.
\item A.~Agrachev, Y.~Sachkov, \emph{Control Theory from the Geometric
Viewpoint}, Springer 2004, Ch.~20--21.
\item R.~Vinter, \emph{Optimal Control}, Birkh\"auser 2000 (endpoint
constraints, strong local optimality).
\item \texttt{costate\_common/ms\_conjugate\_test.m}; \texttt{DRO\_tulip/reviews/huber\_review\_adjudicated\_2026-09-06.md};
\texttt{doc/extremal\_and\_local\_min\_survey.md}.
\end{itemize}
\end{document}
